\documentclass[11pt,a4paper]{article}

\usepackage[utf8]{inputenc}
\usepackage[indonesian]{babel}
\usepackage[a4paper,top=2.2cm,bottom=2.2cm,left=2.2cm,right=2.2cm]{geometry}
\usepackage{graphicx}
\usepackage{booktabs}
\usepackage{xcolor}
\usepackage{tcolorbox}
\usepackage{enumitem}

\definecolor{primaryblue}{RGB}{24, 43, 73}
\definecolor{secondaryblue}{RGB}{41, 128, 185}
\definecolor{goldaccent}{RGB}{211, 84, 0}
\definecolor{boxbg}{RGB}{245, 248, 252}

\begin{document}
\thispagestyle{empty}

\begin{center}
    \vspace*{0.8cm}
    \includegraphics[width=3.0cm]{../assets/logo_unib.png}\\[0.6cm]
    {\large \textbf{\color{secondaryblue}UNIVERSITAS BENGKULU -- FAKULTAS TEKNIK}}\\
    {\large \textbf{\color{secondaryblue}PROGRAM STUDI S1 TEKNIK ELEKTRO}}\\[1.1cm]
    \rule{\textwidth}{1.5pt}\\[0.5cm]
    {\Huge \textbf{\color{primaryblue}KOMPILASI MODUL AJAR OBE}}\\[0.3cm]
    {\LARGE \textbf{\color{goldaccent}ANALISIS DERET WAKTU DAN PERAMALAN}}\\
    {\large \textbf{\color{primaryblue}SISTEM TENAGA LISTRIK \& ENERGI TERBARUKAN (TEE-412)}}\\[0.45cm]
    \rule{\textwidth}{1.5pt}\\[0.9cm]
\end{center}

\begin{tcolorbox}[colback=boxbg,colframe=primaryblue,arc=3mm,title=\textbf{Identitas Karya Ciptaan -- Lampiran Contoh Ciptaan HKI (Buku / Modul Ajar)}]
\begin{tabular}{p{4.2cm}p{10.8cm}}
\textbf{Pencipta} & : Ir. Novalio Daratha, S.T., M.Sc., Ph.D. (NIP. 197811102003121002) \\
\textbf{Institusi} & : Jurusan Teknik Elektro, Fakultas Teknik, Universitas Bengkulu \\
\textbf{Standar Kurikulum} & : Outcome-Based Education (OBE) -- LAM Teknik / IABEE (16 Minggu, 3 SKS) \\
\textbf{Portal Akses Terbuka} & : \texttt{https://www.ndaratha.my.id/deret-waktu} \\
\textbf{Repositori Publik} & : \texttt{https://github.com/novatha/deret-waktu-tee412} \\
\textbf{Tahun Terbit} & : 2026 (Edisi 1 -- Ganjil 2026/2027) \\
\end{tabular}
\end{tcolorbox}

\vspace{0.7cm}
\noindent{\large\textbf{\color{primaryblue}Daftar Isi Kompilasi (Terindeks Bookmark PDF)}}
\vspace{0.2cm}

\begin{enumerate}[label=\textbf{Bagian \Roman*.}, leftmargin=4.2em, itemsep=4pt]
    \item \textbf{Rencana Pembelajaran Semester (RPS OBE 16 Minggu)} -- CPL 1--4, CPMK 1--5, Bobot 100\%, \& 8 Tema Capstone.
    \item \textbf{Buku Panduan Praktikum (12 Bab Komputasi Cerdas)} -- Implementasi Python, Google Colab, \& 3 Dataset CSV Terkurasi.
    \item \textbf{Panduan Proyek Akhir IEEE, Template \texttt{IEEEtran}, Kisi-Kisi UTS/UAS, \& Instrumen Evaluasi OBE-CQI.}
    \item \textbf{13 Modul Slide Perkuliahan Beamer} -- Minggu 1--7 (Dekomposisi s.d. Discrete Kalman Filter) \& Minggu 9--14 (FFT/Wavelet s.d. BESS MPC).
    \item \textbf{13 Set Problem Sets (Taksonomi Bloom C1--C6) \& 13 Lembar Kerja Mahasiswa (Worksheets).}
\end{enumerate}

\vfill
\begin{center}
    {\small\color{gray}Hak Cipta \copyright\ 2026 Ir. Novalio Daratha, S.T., M.Sc., Ph.D. / Universitas Bengkulu.\\
    Disusun untuk Pencatatan Hak Cipta DJKI melalui LPPM Universitas Bengkulu.}
\end{center}

\end{document}
